\documentclass[a4paper,12pt, french, noSujetB]{article}

\usepackage{layout_evaluation}

\classe{Seconde}
\titre{Développement décimal d'un nombre réel}
\evalnumber{1}
\evaltype{Exercice autonome}
\date{septembre 2026}

\begin{document}


\emph{La calculatrice n'est pas nécessaire. On posera les calculs utilisés.}

\vspace{-3mm}

\subsection*{Exercice 1 $\ast$}
\begin{enumerate}
    \item Justifier que les nombres suivants sont des décimaux : $\frac{7}{4}, \frac{2}{5}, \frac{9}{125}, \frac{7741}{500}$.
    \item Justifier que les nombres suivants ne sont pas décimaux : $\frac{1}{3}, \frac{3}{11},  \frac{8}{37}$.
    \item Dans quel ensemble de nombre (le plus petit possible) sont ces nombres ? On donnera le nom de l'ensemble et son symbole mathématique.
    \item Que remarque-t-on ? Formule une conjecture sur le développement décimal de ces nombres.
\end{enumerate}

\subsection*{Exercice 2 $\ast\ast$}
On s'intéresse à la proposition suivante : \guillemotleft Les nombres rationnels qui ne sont pas décimaux ont un développement décimal périodique \guillemotright.

Prenons l'exemple $x = 0.1111\dots$, que l'on note $x = 0.\overline{1}$.
\begin{enumerate}
    \item Calculer $10x$ et $10x-x$.
    \item En déduire un écriture de $x$ sous forme de fraction. Vérifier à l'aide de la calculatrice.
    \item Appliquer une méthode similaire à $y=0.121212\dots=0.\overline{12}$.
    \item Décrire une méthode générale pour trouver le rationnel égal à une écriture décimale périodique.
\end{enumerate}
\end{document}
